\documentclass[10pt,a4paper]{amsart}
\usepackage[margin=19mm]{geometry}
\usepackage{amsmath,amssymb,mathtools}
\usepackage[scr=rsfs,cal=euler]{mathalfa}
\usepackage{xfrac}
\usepackage[unicode,hidelinks,bookmarks=false]{hyperref}
\newcommand{\ie}{i.e.\ }
\newcommand{\cf}{cf.\ }
\newcommand{\resp}{respectively\ }
\newcommand{\Subsection}{\subsection}
\providecommand{\nobreakdash}{\nobreak\hbox{-}}
\setlength{\emergencystretch}{2em}
\multlinegap0pt
\usepackage{bm}
\usepackage{bbm}
\usepackage{dsfont}
\usepackage{xspace}
\usepackage{cancel}
\usepackage{mathtools}
\usepackage{amsthm}
\makeatletter
\@ifundefined{theorem}{\newtheorem{theorem}{Theorem}}{}
\makeatother
\newcommand{\Int}{\operatorname{int}}
\title[Counting subgraphs of coloring graphs using shadow graphs]{Counting subgraphs of coloring graphs using shadow graphs}
\author{Simon MacLean}
\date{Source: arXiv:2505.05807v2 (CC BY 4.0)}
\begin{document}
\maketitle

\noindent\emph{Source and licence.} Counting subgraphs of coloring graphs using shadow graphs. Version arXiv:2505.05807v2 (CC BY 4.0), distributed under the Creative Commons Attribution 4.0 International licence, as recorded in the arXiv licence metadata for this exact version and shown on the abstract page https://arxiv.org/abs/2505.05807v2. Full rights notice: licenses/arxiv-2505.05807-CC-BY-4.0.txt.\par\smallskip

\noindent\textit{Editorial scope.} Counted unit: the Theorem whose source label is \texttt{thm:hypercube-trees} in the article (source0.tex lines 1034--1040, source label \texttt{thm:hypercube-trees}), delivered together with the article's own complete original proof of it (source0.tex lines 1042--1058, one \texttt{proof} environment of 17 lines). No mathematics is written, repaired or re-derived: statement and proof are the article's own text line for line. Required context retained uncounted: none. Every definition, display and label of those lines is retained unchanged. Omitted from this package and not redistributed: 1--1033, 1041--1041, 1059--1403. Nothing inside a retained excerpt is omitted or altered: every retained body is delivered exactly as the article's lines are written, so no omission receipt of any kind accompanies this package. Citations: the retained excerpts carry no citation command at all, so no bibliography entry of the article is transcribed and no reference is redistributed. Reference adaptation: none was needed and none was made; every reference inside the delivered unit resolves inside it, so no unresolved reference or citation is produced. Printed slips kept literal: no printed word, symbol, hypothesis or number of the counted statement or of its proof was altered. Layout and class: the article's own class is \texttt{amsart} with packages including fullpage, times, colonequals, amsmath, amssymb, amsthm, url, xcolor, inputenc, babel, amsrefs, bbm, enumerate, tikz; the wrapper is the lane's pinned \texttt{amsart} editorial wrapper at 10pt on A4, which restates the theorem-like environments the retained text needs and adds the article's own macro definitions that the retained text needs (\texttt{\textbackslash Int}), with extra wrapper packages (bm, bbm, dsfont, xspace, cancel, mathtools). The delivered numbering is the unit's own. Assets: no retained span contains a figure environment, a graphics-inclusion command, a PDF-page inclusion, a table or a TikZ picture; no image file is copied into this package, the package declares no asset dependency and no image file is redistributed with it. Licence: version arXiv:2505.05807v2 of this article is distributed under the Creative Commons Attribution 4.0 International licence, as recorded in the arXiv licence metadata for this exact version and shown on the abstract page https://arxiv.org/abs/2505.05807v2. A full rights notice is delivered at licenses/arxiv-2505.05807-CC-BY-4.0.txt. Characters: every retained excerpt is pure ASCII, so the delivered engine, XeLaTeX with its default Latin Modern text font, reports no missing character.
\begin{theorem}\label{thm:hypercube-trees}
The $d$-dimensional hypercube polynomial for a tree $T$ is
$$
\pi_T^{(Q_d)}(k) =\frac{1}{2^d} \sum_{S \subseteq V(T), |S| = d} k (k - 1)^{n + d - 1 - \deg(S)} (k - 2)^{\deg(S) - \Int(S)} (k - 3)^{\Int(S)},
$$
where $\deg(S) = \sum_{v \in S} \deg(v)$ and $\Int(S)$ is the number of interior edges, those with two endpoints in $S$.
\end{theorem}

\begin{proof}
For each set $S=\{v_1,v_2,\dots,v_d\}$ of vertices in $T$, we:
\begin{itemize}
    \item Replace each vertex $v_i$ in $S$ with two shadow vertices $v_{i,1}$ and $v_{i,2}$.
    \item Add edges $v_{i,1}v_{i,2}$ for all $i\le d$, as well as all edges in $\{v_{i,1},v_{i,2}\}\times N(v_i)$.
    \item Add edges $\{v_{i,1},v_{i,2}\}\times \{v_{j,1},v_{j,2}\}$ for all $v_iv_j\in E(T)$.
\end{itemize}

To color this graph, we begin by assigning any of the $k$ colors to $v_{1,1}$, and any of the $k-1$ remaining colors to $v_{1,2}$. Then, traversing the tree, we make the following coloring choices:
\begin{itemize}
    \item $k-2$ choices for any vertex not in $S$ adjacent to a shadow $K_2$ that has already been colored.
    \item $(k-1)(k-2)$ choices for both vertices in a shadow $K_2$ which is not adjacent to another shadow $K_2$ that has already been colored.
    \item $(k-2)(k-3)$ choices for both vertices in a shadow $K_2$ which is adjacent to another shadow $K_2$ that has already been colored.
    \item $k-1$ choices for a vertex not in $S$ which is not adjacent to a shadow $K_2$ that has already been colored.
\end{itemize}
In total, there will be a single factor of $k$. There will be one factor of $k-3$ for each shadow $K_2$, minus one for each connected component, as the first pair to be assigned a color in a connected component will have $(k-1)(k-2)$ choices. This is equivalent to the number of edges in $T$ with two endpoints in $S$. There will be a factor of $k-2$ for all vertices not in $S$ that are adjacent to a vertex in $S$, except for the first vertex to receive a color. There will be an additional factor of $k-2$ for all pairs of shadow vertices. This is equivalent to the number of edges in $T$ with at least one endpoint in $S$, which is $\deg(S)-\Int(S)$. The remaining vertices contribute a factor of $k-1$. Dividing by $2^d$ corrects for the over-counting due to permuting the colors in the color palette of each shadow vertex.
\end{proof}

\end{document}
